% bulk_one_defect.tex -- round 400, reviewer 5.4 after the
% r399 circularity verdict: threshold / phase form of the
% frame inequality for P1 = ind_-(A0) <= 1.
% Outcome: FORM_T REFUTED (rank-1) / FORM_P REFUTED.
% Builds with pdflatex.
% Companion machine check: verify_bulk_one_defect.py.
% Discovery census: experiments/tfpt-discovery/bulk_one_defect_probe.py.
% NO RH CLAIM.  Research documentation, not a theorem of RH.
\documentclass[11pt]{article}

\usepackage[a4paper,margin=2.7cm]{geometry}
\usepackage{amsmath,amssymb,amsthm}
\usepackage{microtype}
\usepackage{xcolor}
\usepackage[colorlinks=true,linkcolor=blue!50!black,urlcolor=blue!50!black]{hyperref}

\newtheorem{lemma}{Lemma}
\newtheorem{prop}{Proposition}
\theoremstyle{remark}
\newtheorem{remark}{Remark}

\newcommand{\R}{\mathbb{R}}
\newcommand{\indm}{\operatorname{ind}_{-}}
\newcommand{\half}{\tfrac12}
\newcommand{\tr}{\operatorname{tr}}

\title{\bfseries Bulk one-defect, refuted as rank-one\\[2mm]
\large The threshold frame and the $C_m$-phase screen after
the r399 circularity verdict}
\author{Stefan Hamann}
\date{August 29, 2026}

\begin{document}
\maketitle

\begin{abstract}\noindent
After energy decay of the centred mass difference was refuted
(r399) and the $3\times 3$ edge chart was written exactly (r401),
the remaining bulk statement on the selected sequence is
$\indm(A_0)\le 1$ with $A_0=R_{N-3}-\half I$.
Two honest candidate forms are adjudicated.
Form~T: a source-defined rank-one frame
$\langle x,A_0 x\rangle\ge c\|x\|^2-C|\ell(x)|^2$ with $c>0$,
equivalently $A_0$ a positive matrix minus rank one.
Form~P: a phase statistic of the r399 Dirichlet coefficients
$C_m$ that predicts $\indm$ rather than the energy $E$.
What is proved: the frame implication, rank-$r$ interlacing,
and Weyl's inertia bound (finite linear algebra).
What is refuted: both candidate forms as routes to P1.
The Bernstein--Szeg\H{o} dual on the holes already has
$\sim 49$ negatives; the prime source lifts $48$ of them
(a high-rank positive repair, numerical rank $49$ on that
subspace).  Rank one can kill at most one negative.
The genuine negative mode of $A_0$, when present, is
orthogonal to the centred difference (cosine $0.0075$).
Consecutive $C_m$-signs do not separate $\indm=1$ from
$\indm=21$.  Dead $\chi$ satisfy P1: they die at
$\mathrm{sch}>0$ (r401), not in the bulk.
No RH claim.
\end{abstract}

\medskip\noindent
\textbf{Status.}\enspace
Lemmas~\ref{lem:frame}--\ref{lem:weyl} are proved (finite
identities).
Propositions~\ref{prop:rank} and~\ref{prop:phase} are finite
machine censuses.
No RH claim.  No $L^*$ claim.  No $R^\dagger\succ\half I$ claim.

\section{The object}\label{sec:obj}

Write $R=R_{N-3}$ for the dual Christoffel--Darboux Gram of
the first $N_w-3$ dual OP columns on the hole nodes $Y$, and
$A_0=R-\half I$.
P1 is $\indm A_0\le 1$, equivalently $\lambda_2(R)\ge\half$.
Round 398 found a macroscopic cluster of $R$ at $\half$ from
above (FRAME-A: $36$ eigenvalues in $[\half,3/5]$).
Round 399 gave an exact Dirichlet representation of the
centred difference $d$ after Fej\'er, with energy $E$ at
quadratic-mean, growing, and with any bound forcing $E\to 0$
typed RH-near.
The reviewer remainder after that circularity verdict is
whether P1 still follows from a \emph{threshold} (not a
decay) or from the \emph{phases} of the coefficients $C_m$.

The Bernstein--Szeg\H{o} reference $R^{\mathrm{ref}}$ is the
same dual CD Gram built from full-grid Fej\'er weights on the
cosine lattice (r390: $\gamma_k=1/4$ at residual $2\cdot 10^{-6}$),
restricted to the same $Y$.
The perturbation is $B=R-R^{\mathrm{ref}}$.
(An occupied-Fej\'er dual with a tiny floor on $\nu$-nodes
saturates $R^{\mathrm{ref}}=I$ because the dual chain depth
$N_w-3$ exceeds $|Y|$; that degeneracy is disclosed and is
\emph{not} the r390 object.)

\section{Finite algebra}\label{sec:alg}

\begin{lemma}[frame implies P1]\label{lem:frame}
Let $A$ be real symmetric $n\times n$.
If there exist $c>0$, $C\in\R$ and a linear form $\ell$ such that
\[
\langle x,Ax\rangle
\;\ge\;
c\,\|x\|^2-C\,|\ell(x)|^2
\qquad\text{for all }x,
\]
equivalently $A+C\,\ell\otimes\ell\succeq c I$, then $\indm A\le 1$.
\end{lemma}

\begin{proof}
$A+C\,\ell\otimes\ell$ is positive definite, so $A$ differs
from a positive-definite matrix by a rank-one update.
Cauchy interlacing (or Haynsworth on a $1\times 1$ border)
permits at most one negative eigenvalue.
\end{proof}

\begin{lemma}[rank-$r$ interlacing]\label{lem:interlace}
A symmetric rank-$r$ update changes $\indm$ by at most $r$.
In particular a rank-one update kills or creates at most one
negative eigenvalue.
\end{lemma}

Over $\mathbb{Q}$: $\mathrm{diag}(-1,-1)+e_0 e_0^T=\mathrm{diag}(0,-1)$
kills exactly one negative.

\begin{lemma}[Weyl]\label{lem:weyl}
$\indm(A+B)\le\indm A+\indm B$.
\end{lemma}

On FRAME-A the Weyl bound through $R^{\mathrm{ref}}$ is $49+7=56$,
useless as a proof of P1.

\section{Form T, refuted as rank one}\label{sec:T}

\begin{prop}[rank anatomy]\label{prop:rank}
On FRAME-A the Bernstein--Szeg\H{o} dual on $Y$ has
$\indm(R^{\mathrm{ref}}-\half I)=49$.
MAIN kills $48$ of those negatives ($\indm A_0=1$).
The lift $M=P^{\mathrm{ref}}_- B P^{\mathrm{ref}}_-$ on the
$49$-dimensional negative space of the reference has
effective rank $r_{\mathrm{eff}}=\|M\|_F^2/\|M\|_{\mathrm{op}}^2=9.55$,
leading-share $0.105$, and numerical rank $49$.
The spectral-free surrogate (whole $B$) has $r_{\mathrm{eff}}=14.91$.
Core-$42$ has $\indm A_0\in\{0,1\}$ on $42/42$ ($28$ P1 and
$14$ vacuous), with $\lambda_2(R)-\half$ in
$[4.3\cdot 10^{-8},2.0\cdot 10^{-4}]$.
\end{prop}

Lemma~\ref{lem:interlace} then forbids a rank-one $B$:
killing $48$ negatives requires rank at least $48$.
The $36$-mode $\half$-cluster of MAIN is the remnant of that
high-rank lift, sitting razor-thin above $\half$.
A source-defined $\ell$ built from the centred difference
sampled on $Y$ is orthogonal to the genuine negative mode
(cosine $0.0075$); the frame of Lemma~\ref{lem:frame} fails
for every $C$ in a sealed scan, while the tautological
eigenvector works.
The distinguished pair (r367) carries only mass $0.151$ of
that mode: the defect is rest-hosted, not an edge functional
on two nodes.

The exact threshold that implies P1, once
$R^{\mathrm{ref}}$ is \emph{not} already P1, is not a bound
on a sub-band energy of $C_m$.
Low-band energy of $C_m$ is already $\sim 0$ on MAIN, and
$d$ is orthogonal to $v_-$, so Montgomery--Vaughan on short
bands---unconditional, and sharp on short intervals---does
not see P1.
Proving $\lambda_2(R)\ge\half$ by energy methods is the same
wall as r399, now with a named orthogonality.

\section{Form P, refuted as a predictor}\label{sec:P}

\begin{prop}[phase screen]\label{prop:phase}
Consecutive-sign correlation of $C_m$ does not track
$\indm$.
FRAME-A has $\rho_1=-0.341$ and $\indm=1$; scramble
(seed $1$) has $\rho_1=-0.363$ (more anti-correlated) and
$\indm=21$; energies $1.81$ versus $3.39$ (factor $<2$).
Dead $\chi_3$-$15$ has $\rho_1=-0.333$ and $\indm=1$.
Flip rate, lag-$2$, HHI and dyadic band shares likewise
fail to separate the two inertias.
\end{prop}

A source-side bound on Weyl phases of
$\sum\Lambda(n)n^{-1/2-it}$ is therefore not a route to P1:
even a theorem-grade phase law would control the wrong
statistic.
The honesty gate is not reached, because the empirical
predictor is already dead.

\section{Named kills}\label{sec:kills}

Scramble: $\indm=21$, reference negatives $35$, only $14$
killed, $r_{\min}=0.0036$ (the cluster punches through
$\half$).  The perturbation is not rank one and does not
finish the lift.
Two-period ($S=21$): $\indm=4$, killed $0$ (no
Bernstein--Szeg\H{o} lift).
Mutants: omit $\ell$ (Lemma~\ref{lem:frame} hypothesis
fails); threshold $\times 2$ asks $\lambda_2(R)\ge 1$, false
at $\lambda_2=0.50004$.
Dead $\chi_3$ (five rungs) and dead $\chi_4$-$20$:
$\indm A_0\in\{0,1\}$ on all of them.
Living $\chi_3$ $37/37$ and living $\chi_4$ $41/41$ likewise.
P1 is not world-separating.
The death of dead $\chi$ sits at $\mathrm{sch}>0$ in the
r401 chart, i.e.\ entirely in the edge.

\section*{Machine check}

Standalone: Fractions frame, omit-$\ell$, rank-one
interlacing, Weyl, threshold $\times 2$.
Construction: FRAME-A lift rank, orthogonality, source-frame
failure, scramble $21$, two-period, dead/living $\chi$,
phase-blindness.
Discovery census: \texttt{bulk\_one\_defect\_probe.py}
(\texttt{--smoke} / full), core-$42$, Twin, living and dead
$\chi$.
Exit line: \texttt{BULK ONE DEFECT VERIFIED}.

\section*{Claim boundary}

Finite identities for the frame implication, rank-$r$
interlacing and Weyl; a named census that the bulk
perturbation is a high-rank lift of a Bernstein--Szeg\H{o}
dual, not a rank-one defect, and that $C_m$-phases do not
predict $\indm$; named scramble / two-period / mutant /
dead-$\chi$ controls.
No statement about zeros of $\zeta(s)$, in either direction.
No RH claim.

\begin{thebibliography}{9}
\bibitem{r367}
S.~Hamann,
\textit{Final two-rank inertia},
sealed probe, August 2026.
\bibitem{r390}
S.~Hamann,
\textit{$G_\varepsilon^\mu$ / Bernstein--Szeg\H{o} Fej\'er},
sealed probe, August 2026.
\bibitem{r398}
S.~Hamann,
\textit{High-moment inertia},
standalone note, August 2026,
\texttt{rh/problem/high\_moment\_inertia.tex}.
\bibitem{r399}
S.~Hamann,
\textit{Source Weyl energy},
standalone note, August 2026,
\texttt{rh/problem/source\_weyl\_energy.tex}.
\bibitem{r401}
S.~Hamann,
\textit{The augmented edge signature},
standalone note, August 2026,
\texttt{rh/problem/edge\_signature.tex}.
\end{thebibliography}

\end{document}
